\scanpage{34}
$ {}_nT_0$, on peut le prolonger en un sous-groupe \inserted{fini} étale $H(n)$ de $G$, tel que $H(m)={}_{\unsure{m|n}}H(n)$ si $n\mid m$. \margin{$G$ lisse$/k$}La famille des sous-groupes fermés lisses $\underline{\mathrm{Centr}}_G(H(n))$ de $G$ est stationnaire, sa valeur limite $\deleted{D}C$ est telle que $D_0=\mathrm{Centr}_{G_0}(T_0)=C_0$. De plus, on voit que $C_1$ a $\overline{\mathrm m}$ rang réductif et nilpotent que $G_1$ [i.e. \deleted{\unsure{contient}} $C_{\bar{s}_1}^{\deleted{\unsure{\circ}}}$ contient un Cartan de $G_{\bar{s}_1}$] : c'est un \uline{lemme à rapporter}.

\margin{$\ast$}
Il s'ensuit que \deleted{$\rho_n^{\unsure{1}}(G_1)=\rho_n(\unsure{C_1}),\ \rho_r(G_1)=\rho_r(C_1)$} $\rho_r,\rho_u,\rho_{ab},\rho_n,\rho_{na}$ sont les $\overline{\mathrm m}\mathrm{s}$ pour $G_1$ et $C_1$
\[
\begin{aligned}
\deleted{\rho_{\unsure{u|n}}(G_1)=\rho_n(C_1)}\\
\deleted{\rho_u(G_1)=\rho_u(C_1)}\\
\deleted{\rho_{ab}(G_1)=\rho_{ab}(C_1)}\\
\deleted{\rho_r(G_1)=\rho_r(C_1)}
\end{aligned}
\]
et sont aussi les $\overline{\mathrm m}\mathrm{s}$ pour $C_0$ et $G_0$, donc pour les formules 1) à 5), on est ramené au cas où $G=C$, i.e. $G_0$ \deleted{à fibres} est nilpotent. Alors la relation $\boxed{\rho_n\geq\rho'_n}$ \unsure{triviale}, car $\rho_n=d=d'$

\uline{Lemme} $C$ gpe nilpotent connexe lisse sur corps alg. clos $k$. \margin{\unsure{peut|pour} s'énoncer sans alg. clos\ldots}
% slanted
Alors $C/\overbrace{\underline{\mathrm{Centre}}(C)_{\mathrm{réd}}^\circ}^{\unsure{Z=|Z}}$ est unipotent. Pour tt gpe de type multiplicatif $H$, $\underline{\mathrm{Hom}}_{k\text{-gr}}(H,C)\simeq\underline{\mathrm{Hom}}_{k\text{-gr}}(H,Z)$. En particulier, pour $n$ premier à la car, $ {}_nC\simeq{}_nZ$ est un \unsure{groupe|sous-groupe} comm. étale de $C$, central,
% archive sheet number: 33
